\documentclass[11pt]{article}

\usepackage[top=0.5in, bottom=0.5in, left=0.85in, right=0.7in]{geometry}
\usepackage{authblk}
\usepackage{hyperref}
\usepackage[utf8]{inputenc}
\usepackage{amsmath}
\usepackage{amsfonts}
\usepackage{amssymb}
\usepackage{siunitx}
\usepackage{graphicx}
\usepackage{subcaption}
\usepackage{float}
\usepackage{indentfirst}
\usepackage{longtable}
\usepackage{makecell} % Add this in your preamble
\usepackage{needspace}


\usepackage[backend=biber, style=numeric, sorting=unsrt]{biblatex}
\addbibresource{References.bib} % Ensure your .bib file name matches

\usepackage{setspace}
\onehalfspacing
\setlength{\parskip}{1em}
\title{ % \large 
  \textbf{Computational Economics: Problem Set 5}}
  % Boosting Wearable Sales: How Technology Development and Research Funding Impact Consumer Adoption and Market Growth

\author{\\Emma (Sunwoo) Cho}

\makeatletter
\let\inserttitle\@title
\let\insertauthor\@author
\makeatother

\begin{document}

\begin{center}
  \LARGE{\inserttitle}

  \insertauthor
\end{center}
%------------------------------------------------------------------------------
\section*{Part A \& B}
\begin{itemize}
    \item 3 figures and 3 tables, interpretations for each of 3 models (variation of each)
    \item Codes (in R), clean function; model fitting function; visualisation function
\end{itemize}

\section*{Models and Specifications}
\begin{itemize}
    \item The economic problem is looking the impact of the tax cut event in Kansas in 2012 (occurred in Q2) on the ln (state GDP per capita).
    \item Here, the main identification strategy for Augmented Synthetic Control (ASCM) variants to estimate the Average Treatment Effect on the Treated (ATT): the post-2012Q2 average gap between Kansas and its synthetic control.
    \item 3 specifications:
    \begin{enumerate}
        \item ASCM (Ridge): starts from SCM and augments with a ridge outcome model, allowing limited extrapolation to improve pre-trend fit while penalising large deviations.
        \item ASCM (Fixed Effects): SCM with an additive unit+time structure (demeaned/two-way FE augmentation) to remove level differences and align dynamics.
        \item ASCM (Generalised Synthetic Control Method (Gen-SC) - Matrix Completion method): augments with a low-rank interactive fixed-effects model to capture latent factors.
    \end{enumerate}

\section*{Figures}
\begin{itemize}
    \item 3 figures and 3 tables, interpretations for each of 3 models (variation of each)
    \end{itemize}
    
    \begin{figure}
        \centering
        \includegraphics[width=0.8\linewidth]{fig1_ascm_ridge.png}
        %\caption{Caption}
        %\label{fig:placeholder}
    \end{figure}

    \begin{figure}
        \centering
        \includegraphics[width=0.8\linewidth]{fig2_ascm_fixedeff_gap.png}
       % \caption{Caption}
        %\label{fig:placeholder}
    \end{figure}

        \begin{figure}
        \centering
        \includegraphics[width=0.8\linewidth]{fig3_ascm_gsyn.png}
        %\caption{Caption}
        %\label{fig:placeholder}
    \end{figure}
\end{itemize}

\begin{itemize}
    \item figure 1: for re-treatment, fit gaps oscillate around 0 very tightly, which means it has a good pre-trend balance. For post-treatment: the gap drops and stays negative approximately (−0.03 to −0.06) for most post years. Kansas’s GSP per capita is about 3–6\% lower than its synthetic control after the 2012 tax cuts, with bands mostly below zero. This is consistent with a negative effect.
    \item figure 2: showing very similar as figure 1; it is a robustness check form using a fixed effect augmentation and it gives very similar negative impacts, supporting the ridge result.
    \item figure 3: generalised synthetic control using matrix completion approach - for pre-treatment, it balanced on average, but with a bit more wiggle in late pre-period. For post-treatment, the line shows a small positive bump early, then meanders; the gray band is wide and often covers zero. 
\end{itemize}

\begin{table}[t]
\centering
\caption{Average Treatment on Treated (ATT) estimates across 3 specifications (model fitting results)}
\label{tab:ascm_att_simple}
\begin{tabular}{lccc}
\hline
Model & N\_post & ATT (log) & Percent \\
\hline
ASCM (Ridge)         & 16 & 0.036 & 3.7 \\
ASCM (Fixed Effects) & 16 & 0.064 & 6.6 \\
ASCM (GSYN)          & 16 & 0.094 & 9.9 \\
\hline
\end{tabular}
\end{table}

\subsection*{Interpretation}
The table shows the average treatment effects on the log scale and convert them to percentages via
$100\,[\exp(\text{ATT})-1]$,
And these estimates are in the percentage column.

\begin{itemize}
  \item \textbf{ASCM (Ridge):} $\text{ATT}=0.036 \; \approx 3.7\%$ higher outcome for treated units in the post period.
  \item \textbf{ASCM (Fixed Effects):} $\text{ATT}=0.064 \; \approx 6.6\%$.
  \item \textbf{ASCM (GSYN):} $\text{ATT}=0.094 \; \approx 9.9\%$.
\end{itemize}

\noindent All three estimates are positive and increase in magnitude from Ridge to Fixed Effects one to generalised synthetic control method. This pattern is consistent with a robust positive effect and with larger estimated effect after allowing richer latent trends (generalised SC accommodates interactive factors) compared to stronger shrinkage (ridge tends to be conservative).

\clearpage
\Needspace{0.92\textheight} % keep both figures on the same page

\end{document}